\clearpage
\VerdierPhysicalPageAnchor{0269}%
\begingroup
\small
\setlength{\parindent}{0pt}%
\setlength{\parskip}{0pt}%
\setlength{\emergencystretch}{3em}%
\newcommand{\VerdierContentsRow}[3]{%
  \par\noindent
  \hangindent=#1\relax\hangafter=1
  \hspace*{#1}#2\nobreak
  \leaders\hbox{\kern.32em.\kern.32em}\hfill
  \nobreak\hbox{\kern.65em #3}\par}
\newcommand{\VerdierContentsChapter}[1]{%
  \par\addvspace{.55\baselineskip}\noindent\textbf{#1}\par\smallskip}
\newcommand{\VerdierContentsSection}[2]{%
  \par\addvspace{.45\baselineskip}\noindent
  \hfill\textbf{#1.}\enspace #2\hfill\mbox{}\par\smallskip}

\VerdierSemanticAnchor{toc-ch2}%
\VerdierContentsChapter{Chapter II: Triangulated Categories.}
\VerdierSemanticAnchor{toc-ch2-sec1}%
\VerdierContentsSection{1}{Definitions and First Properties.}
\VerdierSemanticAnchor{toc-ch2-sec1-1}%
\VerdierContentsRow{.5em}{\textbf{1.1.}\enspace Definition of Triangulated Categories.}{93}
\VerdierSemanticAnchor{toc-ch2-sec1-2}%
\VerdierContentsRow{.5em}{\textbf{1.2.}\enspace First Properties of Triangulated Categories.}{97}
\VerdierSemanticAnchor{toc-ch2-sec1-3}%
\VerdierContentsRow{.5em}{\textbf{1.3.}\enspace Examples of triangulated categories.}{106}

\VerdierSemanticAnchor{toc-ch2-sec2}%
\VerdierContentsSection{2}{Localization in triangulated categories.}
\VerdierSemanticAnchor{toc-ch2-sec2-1}%
\VerdierContentsRow{.5em}{\textbf{2.1.}\enspace Multiplicative systems of morphisms.}{111}
\VerdierSemanticAnchor{toc-ch2-sec2-2}%
\VerdierContentsRow{.5em}{\textbf{2.2.}\enspace Construction of the localized category.}{117}
\VerdierSemanticAnchor{toc-ch2-sec2-3}%
\VerdierContentsRow{.5em}{\textbf{2.3.}\enspace Properties of the localization functor.}{125}

\VerdierSemanticAnchor{toc-ch2-sec3}%
\VerdierContentsSection{3}{Abelianization of triangulated categories.}
\VerdierSemanticAnchor{toc-ch2-sec3-1}%
\VerdierContentsRow{.5em}{\textbf{3.1.}\enspace Construction of the category $\mathrm A(\mathcal D)$.}{135}
\VerdierSemanticAnchor{toc-ch2-sec3-2}%
\VerdierContentsRow{.5em}{\textbf{3.2.}\enspace Properties of $\mathrm A(\mathcal D)$.}{137}

\VerdierSemanticAnchor{toc-ch2-sec4}%
\VerdierContentsSection{4}{Spectral objects.}
\VerdierSemanticAnchor{toc-ch2-sec4-1}%
\VerdierContentsRow{.5em}{\textbf{4.1.}\enspace Definition of spectral objects.}{145}
\VerdierSemanticAnchor{toc-ch2-sec4-2}%
\VerdierContentsRow{.5em}{\textbf{4.2.}\enspace Mechanism of spectral sequences.}{147}
\VerdierSemanticAnchor{toc-ch2-sec4-3}%
\VerdierContentsRow{.5em}{\textbf{4.3.}\enspace The terms $E\!\begin{bmatrix}n\\-\infty,p,q,+\infty\end{bmatrix}$.}{153}
\VerdierSemanticAnchor{toc-ch2-sec4-4}%
\VerdierContentsRow{.5em}{\textbf{4.4.}\enspace Convergence.}{155}

\VerdierSemanticAnchor{toc-ch3}%
\VerdierContentsChapter{Chapter III: Derived Categories.}
\VerdierSemanticAnchor{toc-ch3-sec1}%
\VerdierContentsSection{1}{Definitions of Derived Categories.}
\VerdierSemanticAnchor{toc-ch3-sec1-1}%
\VerdierContentsRow{.5em}{\textbf{1.1.}\enspace Notation and Terminology.}{157}
\VerdierSemanticAnchor{toc-ch3-sec1-2}%
\VerdierContentsRow{.5em}{\textbf{1.2.}\enspace The Notable Subcategories of Derived Categories.}{159}
\VerdierSemanticAnchor{toc-ch3-sec1-3}%
\VerdierContentsRow{.5em}{\textbf{1.3.}\enspace The canonical $\delta$-functor.}{164}

\VerdierSemanticAnchor{toc-ch3-sec2}%
\VerdierContentsSection{2}{Resolutions.}
\VerdierSemanticAnchor{toc-ch3-sec2-1}%
\VerdierContentsRow{.5em}{\textbf{2.1.}\enspace Elementary morphisms.}{169}
\VerdierSemanticAnchor{toc-ch3-sec2-2}%
\VerdierContentsRow{.5em}{\textbf{2.2.}\enspace Study of certain sets of objects.}{173}
\VerdierSemanticAnchor{toc-ch3-sec2-3}%
\VerdierContentsRow{.5em}{\textbf{2.3.}\enspace Some lemmas on resolutions.}{177}
\VerdierSemanticAnchor{toc-ch3-sec2-4}%
\VerdierContentsRow{.5em}{\textbf{2.4.}\enspace Applications.}{185}

\VerdierSemanticAnchor{toc-ch3-sec3}%
\VerdierContentsSection{3}{Study of $\operatorname{Ext}$.}
\VerdierSemanticAnchor{toc-ch3-sec3-1}%
\VerdierContentsRow{.5em}{\textbf{3.1.}\enspace Different equivalent definitions.}{191}
\endgroup
